\section{Case Study: An IT/Software-Engineering ISD Prototype}
\label{sec:case-study}

Benchmark evaluation (\cref{sec:experiments}) treats each scenario as an independent run---appropriate for comparing agents, but blind to the two characteristics established in \cref{sec:introduction}: real training programs are refined over many sessions (the core problem), and every domain carries its own documents and conventions (the domain-knowledge requirement). The case study fills that gap with a working application in our own enterprise context: internal IT/software-engineering training for new employees, interns, and junior engineers.

\paragraph{Application architecture and session lifecycle.}
The application wraps the core harness in a session-management layer that lets an alignment graph persist as a long-lived artifact at the level of a training program: \texttt{training\_program\_id} is the durable key of the graph, and an optional \texttt{session\_id} resumes an interrupted session. The first session (\textsl{init}) creates $G^{(1)}$; each subsequent session (\textsl{refinement}) loads the latest version, applies the refinement request, and commits $G^{(v+1)}$, exactly the version chain of \cref{sec:memory}. No new instrumentation is needed: trajectories, cascade events, repair history, and graph version diffs are already recorded by the harness.

\paragraph{Domain instantiation.}
Three representative programs were built: \emph{Java/OOP for junior engineers} (3 days, hybrid), \emph{software testing} (2 days, instructor-led), and a \emph{Git/agile mini-project} (2 days, for new project members). Each subject is packaged as a retrieval agent following the plug-and-play scheme of \cref{sec:agentic-rag}---a corpus (0--4 sample rubrics, assignment and starter-repository conventions, SOLID/code-review checklists, pull-request and sprint-simulation conventions), a system prompt describing the corpus scope, and a standard lookup toolset---and registered in $R$. The full workflow is demonstrated from the start of the chain: a vague request (``we need OOP training for juniors before they join the project'') passes through elicitation (\cref{sec:elicitation}) to an approved design brief; the init session generates $G^{(1)}$ with domain knowledge from the corpus; and refinement sessions then handle the requests that arise after deployment (add an objective, tighten a rubric, change the duration, switch the delivery format).

The case study also serves as the substrate of the RQ2 experiment: eight refinement requests were designed across the three programs, spanning the impact-policy spectrum from the narrowest target type (evaluation) to the broadest (needs context), so that the expected cascade scope varies rather than staying flat; the measurement protocol and the four control conditions are given in \cref{sec:experiments}.
